% ============================================================================
% French UI and structural strings for the One Chemistry Book series.
% Loaded by styles/onechemistry.sty when \booklang is fr.
% ============================================================================

% Theorem / statement environment titles
\newcommand{\omnameDefinition}{Définition}
\newcommand{\omnameTheorem}{Théorème}
\newcommand{\omnameProposition}{Proposition}
\newcommand{\omnameLemma}{Lemme}
\newcommand{\omnameCorollary}{Corollaire}
\newcommand{\omnameMethod}{Méthode}
\newcommand{\omnameExample}{Exemple}
\newcommand{\omnameNotation}{Notation}
\newcommand{\omnameRemark}{Remarque}
\newcommand{\omnameExercise}{Exercice}
\newcommand{\omnameExercises}{Exercices}
\newcommand{\omnameProblem}{Problème}
\newcommand{\omnameProblems}{Problèmes}
\newcommand{\omnameProof}{Démonstration}

% Admitted results and solutions appendix headers
\newcommand{\omadmittedtext}{Admis à ce niveau.}
\newcommand{\omsolutionof}{Solution de}
\newcommand{\ompage}{page}
\newcommand{\omnameGotoSolution}{Solution p.}

% Misc text macros
\newcommand{\st}{\text{ tel que }}

% Cover / title page
\newcommand{\bookmaintitle}{One Chemistry Book}
\newcommand{\booktagline}{Le Livre de Chimie pour les gouverner tous}
\newcommand{\bookdescription}{La chimie de la première année d'école à la fin de la licence}
\newcommand{\bookcontributors}{Les contributeurs du One Chemistry Book}
\newcommand{\bookversionlabel}{Version}

% Colophon
\newcommand{\omcolophonsource}{Code source, signalements et contributions~:}
\newcommand{\omcolophonlicense}{Sauf mention contraire, ce livre est placé sous licence Creative Commons CC BY-NC-SA 4.0~:}
\newcommand{\omcolophoncredit}{Merci de créditer ainsi~: One Chemistry Book, One Course (one-course.com).}

% Chemistry boxes (styles/onechemistry.sty): the recall box that opens a
% chapter building on another one, the lab box, the history box, the safety
% box. Placeholder wording by the coordinator (2026-10-04); the Book 1
% `fr` agent owns it and settles the register (the recall title addresses
% the reader in the same register as the book's exercise stems).
\newcommand{\omnameRecall}{Rappel}
\newcommand{\omnameInTheLab}{Au laboratoire}
\newcommand{\omnameHistory}{Histoire des sciences}
\newcommand{\omnameSafety}{Sécurité}

% Solutions appendix part title
\newcommand{\omsolutionspart}{Solutions des exercices}

% Part titles (Années 1–12; parallel to English Grade N)
\@namedef{omstr@part.grade1}{Année 1 (6--7 ans)}
\@namedef{omstr@part.grade2}{Année 2 (7--8 ans)}
\@namedef{omstr@part.grade3}{Année 3 (8--9 ans)}
\@namedef{omstr@part.grade4}{Année 4 (9--10 ans)}
\@namedef{omstr@part.grade5}{Année 5 (10--11 ans)}
\@namedef{omstr@part.grade6}{Année 6 (11--12 ans)}
\@namedef{omstr@part.grade7}{Année 7 (12--13 ans)}
\@namedef{omstr@part.grade8}{Année 8 (13--14 ans)}
\@namedef{omstr@part.grade9}{Année 9 (14--15 ans)}
\@namedef{omstr@part.grade10}{Année 10 (15--16 ans)}
\@namedef{omstr@part.grade11}{Année 11 (16--17 ans)}
\@namedef{omstr@part.grade12}{Année 12 (17--18 ans)}

% Part titles (Licence 1–3; parallel to English Bachelor Year N).
% \textsuperscript, not babel's \iere/\ieme: babel may be absent (see onechemistry.sty).
\@namedef{omstr@part.bachelor1}{Licence 1\texorpdfstring{\textsuperscript{re}}{re} année (18--19 ans)}
\@namedef{omstr@part.bachelor2}{Licence 2\texorpdfstring{\textsuperscript{e}}{e} année (19--20 ans)}
\@namedef{omstr@part.bachelor3}{Licence 3\texorpdfstring{\textsuperscript{e}}{e} année (20--21 ans)}

% Plural names + \omnameChapter, used by the \crefname declarations in
% onechemistry.sty. Without these, cleveref falls back to its English built-ins and
% a French book says "Theorem 3.1" / "Chapter 5" in the middle of French prose.
\newcommand{\omnameDefinitions}{Définitions}
\newcommand{\omnameTheorems}{Théorèmes}
\newcommand{\omnamePropositions}{Propositions}
\newcommand{\omnameLemmas}{Lemmes}
\newcommand{\omnameCorollarys}{Corollaires}
\newcommand{\omnameMethods}{Méthodes}
\newcommand{\omnameExamples}{Exemples}
\newcommand{\omnameNotations}{Notations}
\newcommand{\omnameRemarks}{Remarques}
\newcommand{\omnameChapter}{Chapitre}
\newcommand{\omnameChapters}{Chapitres}
% Section, for \cref{sec:…} (2026-10-04: it had no \crefname, so cleveref
% printed its English built-in "Section" in every edition).
\newcommand{\omnameSection}{Section}
\newcommand{\omnameSections}{Sections}

% siunitx and cleveref keep their English conjunctions unless told otherwise:
% \qtyrange printed "5 to 4186 Hz" and \cref{a,b} printed "Chapitres 25 and 29"
% in the middle of French prose.
% cleveref installs its conjunctions at \begin{document}, so redefine them there.
% siunitx typesets range phrases and list separators in MATH MODE, where a
% bare string loses its spaces (\qtyrange inside $...$ printed "1a100") and a
% non-ASCII character is worse: "a" expands to \`a, which LaTeX reports as
% "Command \` invalid in math mode" and drops from the page. siunitx's own
% defaults are wrapped in \text{} for exactly this reason; these overrides
% must be too. Found on Biology Book 3 fr, 2026-09-05, where the book's 56
% \qtyrange sites made it visible; it was latent in every language edition of
% every book in this repo.
\sisetup{range-phrase={\text{ à }}, list-pair-separator={\text{ et }},
         list-final-separator={\text{ et }}}
\AtBeginDocument{%
  \renewcommand{\crefpairconjunction}{ et }%
  \renewcommand{\creflastconjunction}{ et }%
  \renewcommand{\crefrangeconjunction}{ à }%
  % a \cref spanning two label types builds a group list, which has its
  % own pair of conjunctions: "Théorèmes 2.14 et 2.18 and Lemme 2.17".
  \renewcommand{\crefpairgroupconjunction}{ et }%
  \renewcommand{\creflastgroupconjunction}{ et }%
}

% \today is English without babel (which is optional here), so the cover date
% read "July 27, 2026" on a French book. French writes it as
% "27 juillet 2026", month lowercase.
% (the first of the month is "1\textsuperscript{er}", every other day a bare
% numeral).
\renewcommand{\today}{\ifnum\day=1 1\textsuperscript{er}\else\number\day\fi
  \space\ifcase\month\or janvier\or février\or
  mars\or avril\or mai\or juin\or juillet\or août\or septembre\or octobre\or
  novembre\or décembre\fi\space\number\year}

% LaTeX's own structural names. babel would normally localise these, but it is
% optional here (see onechemistry.sty) — so without these a French book prints
% "Chapter 1" above every chapter title and "Appendix A" over the solutions.
\renewcommand{\chaptername}{Chapitre}
\renewcommand{\partname}{Partie}
\renewcommand{\contentsname}{Table des matières}
\renewcommand{\appendixname}{Annexe}
\renewcommand{\indexname}{Index}

% Periodic-table legend (\omperiodictable[legend], styles/onechemistry.sty).
% Coordinator wording, 2026-10-04 (it used to be hard-coded English in the
% style file, so no edition could translate it); the Book 1 agent of this
% language checks it against the book's own terminology.
\newcommand{\omnamePTmetal}{métaux}
\newcommand{\omnamePTmetalloid}{métalloïdes}
\newcommand{\omnamePTnonmetal}{non-métaux}
\newcommand{\omnamePTs}{bloc s}
\newcommand{\omnamePTp}{bloc p}
\newcommand{\omnamePTd}{bloc d}
\newcommand{\omnamePTf}{bloc f}
\newcommand{\omnamePTalkali}{métaux alcalins}
\newcommand{\omnamePTalkaline}{métaux alcalino-terreux}
\newcommand{\omnamePThalogen}{halogènes}
\newcommand{\omnamePTnoble}{gaz nobles}
